\section{Standards, Frameworks, and Regulatory Governance}
\label{sec:standards}

The transition of LLMs from experimental prototypes to mission-critical agentic deployments has prompted major cybersecurity standards bodies and regulatory agencies to establish governance frameworks~\cite{hasselwander2026agentic}. Harmonizing these frameworks with technical attack surfaces is essential for engineering compliant, defensible systems. Table~\ref{tab:standards_mapping} maps agentic attack surfaces across primary international security standards.

\begin{table}[t]
\centering
\caption{Cross-Framework Harmonization: Mapping Agentic Attack Surfaces to OWASP Top 10 (2025), MITRE ATLAS, and NIST AI RMF.}
\label{tab:standards_mapping}
\footnotesize
\begin{tabularx}{\columnwidth}{p{2.2cm} X X X X}
\toprule
\textbf{Agentic Attack Surface} & \textbf{OWASP LLM 2025~\cite{owasp2025top10}} & \textbf{MITRE ATLAS~\cite{mitre2026atlas}} & \textbf{NIST AI RMF (100-2e2025)~\cite{vassilev2025adversarial}} & \textbf{EU AI Act Compliance Focus} \\
\midrule
\textbf{Direct \& Indirect Injection} & LLM01: Prompt Injection & AML.T0051.000 / .001: Prompt Injection & Evasion: Indirect Context Poisoning & Art. 15: Cyber Resilience \& Robustness \\
\addlinespace
\textbf{RAG \& Knowledge Corruption} & LLM04: Poisoning; LLM08: Embedding Weakness & AML.T0070: Data Poisoning: Retrieval & Poisoning: Training \& Ingestion Stores & Art. 10: Data Governance \& Quality \\
\addlinespace
\textbf{Tool Misuse \& Parameter Injection} & LLM05: Output Handling; LLM06: Excessive Agency & AML.T0053: LLM Plugin Compromise & Integrity: System-Level Subversion & Art. 14: Human Oversight \& Control \\
\addlinespace
\textbf{API \& Cloud Credential Abuse} & LLM06: Excessive Agency; LLM02: Sensitive Info & AML.T0086: Lateral Movement & Unauthorized Privilege Traversal & Art. 15: Access Boundaries \& Isolation \\
\addlinespace
\textbf{Memory \& Backdoor Persistence} & LLM04: Poisoning; LLM07: System Prompt Leak & AML.T0080.000: Conversation Hijacking & Persistence: State Subversion & Art. 12: Record-Keeping \& Audit Logs \\
\addlinespace
\textbf{Multi-Agent Cascades \& Swarms} & LLM01: Prompt Injection; LLM06: Excessive Agency & AML.T0051.001: Indirect Prompt Injection & Distributed Agent Compromise & Art. 14: Systemic Risk Assessment \\
\addlinespace
\textbf{Denial of Cash (DoC) \& Loops} & LLM10: Unbounded Consumption & AML.T0034: Compute Hijacking & Availability: Resource Depletion & Art. 9: Risk Management Systems \\
\bottomrule
\end{tabularx}
\end{table}


\subsection{OWASP Top 10 for LLM Applications (2025 Edition)}
\label{subsec:owasp_mapping}

The 2025 release of the OWASP Top 10 for LLM Applications~\cite{owasp2025top10, promptfoo2025owasp} explicitly reflects the risks introduced by agentic orchestration:
\begin{itemize}
    \item \textbf{LLM01: Prompt Injection:} Direct and indirect injections are recognized as the primary root-cause vulnerability enabling tool hijacking.
    \item \textbf{LLM06: Excessive Agency:} Identified as the primary architectural flaw in agentic deployments, occurring when agents are granted excessive permissions, excessive autonomy, or unconstrained tool access without strict least-privilege scoping.
    \item \textbf{LLM08: Vector and Embedding Weaknesses:} A dedicated entry addressing adversarial manipulation of embeddings, semantic cache poisoning, and access-control failures in vector retrieval.
    \item \textbf{LLM10: Unbounded Consumption:} Addresses Denial of Cash (DoC) and resource exhaustion attacks where loops or high-frequency tool invocations overwhelm computational budgets.
\end{itemize}

\subsection{MITRE ATLAS Matrix Harmonization}
\label{subsec:mitre_atlas}

The MITRE Adversarial Threat Landscape for Artificial-Intelligence Systems (ATLAS)~\cite{mitre2026atlas} catalogs real-world adversary tactics, techniques, and procedures (TTPs). In agentic AI, ATLAS techniques map directly to multi-stage exploit pathways:
\begin{itemize}
    \item \textit{Initial Access \& Execution:} AML.T0051 (LLM Prompt Injection) and AML.T0053 (LLM Plugin / Tool Compromise).
    \item \textit{Persistence:} AML.T0080.000 (LLM Conversation / Memory Hijacking) and AML.T0070 (Data Poisoning of Knowledge Retrieval Stores).
    \item \textit{Lateral Movement:} AML.T0086 (Lateral Movement via API and Cloud Service Connectors).
    \item \textit{Impact:} AML.T0034 (Compute and Resource Hijacking).
\end{itemize}

\subsection{NIST AI Risk Management Framework and International Regulation}
\label{subsec:nist_eu_act}

The National Institute of Standards and Technology (NIST) updated report on Adversarial Machine Learning (NIST AI 100-2e2025)~\cite{vassilev2025adversarial} formalizes trustworthiness dimensions across validity, reliability, safety, and security. For autonomous agents, NIST mandates strict isolation of untrusted data feeds, comprehensive provenance logging, and verifiable boundary controls.

Concurrently, the European Union's \textit{Artificial Intelligence Act (EU AI Act)} establishes legally binding mandates for high-risk autonomous AI systems:
\begin{itemize}
    \item \textbf{Article 14 (Human Oversight):} Mandates that autonomous agents possess built-in operational stops, verifiable confirmation interfaces, and capabilities for human operators to override or halt execution without state corruption.
    \item \textbf{Article 15 (Cybersecurity and Robustness):} Requires autonomous systems to demonstrate resilience against adversarial exploitation, including data poisoning, indirect prompt injection, and unauthorized lateral movement across enterprise networks.
\end{itemize}
Compliance with these mandates necessitates bridging high-level governance policies with the technical defense-in-depth mechanisms detailed in Section~\ref{sec:defenses}.
